\section{Source audit and proposed corrections}
\label{audit:sec}

\subsection{The Gaussian candidates remain conjectures}

The manuscript distinguishes proved reductions from frozen
integer-relation candidates. That distinction is preserved here.
The current $S_6$ statement is
\texttt{cycloquot:conj:S6} in
\texttt{chapters/04-cyclotomic-quotients.tex}; the current $S_8$
statement is \texttt{s8new:conj:S8} in
\texttt{chapters/04-S8-candidate.tex} \cite{Repo}.
The latter is a different vector from an earlier rejected $S_8$
candidate. An objection to that rejected vector is not an objection
to the current one.

The supplied rigorous proximity certificates for the current frozen
vectors have residual intervals containing zero. In particular, a
bound $|R|<10^{-355}$ proves a very strong approximation but does not
prove $R=0$. Repeating floating-point integer-relation searches would
not close that logical gap. Nothing proved in this article supplies
the missing exact cyclotomic relation, so neither conjecture is
promoted to a theorem.

The finite separating functional already constructed for a restricted
level-four quotient is useful but has a limited conclusion. It rules
out a derivation inside the particular stated depth-one-product
shuffle/stuffle presentation. It does not prove that the target
constant lies outside every larger cyclotomic algebra, and it does
not rule out octahedral functional equations or additional
regularized relations. The exact proof of $S_4$ in the manuscript,
with its explicitly verified additional relation families, is a
model for the type of enlargement that would be needed.

\subsection{A convergent depth-three bridge for the Gaussian $S_6$ target}
\label{audit:s6-bridge}

The restricted quotient can be tested against an explicit relation just
outside its defining schemas. This gives a concrete exact reduction to
include in a future certificate, while leaving the $S_6$ conjecture open.
Here the manuscript's harmonic sum and Dirichlet beta function are
\[
 S_p=\sum_{n\ge0}\frac{(-1)^nH_n}{(2n+1)^p}\quad(p>1),
 \qquad
 \beta(s)=\sum_{n\ge0}\frac{(-1)^n}{(2n+1)^s}.
\]
Use the series convention
\[
 \Li_{s_1,\ldots,s_d}(x_1,\ldots,x_d)
 =\sum_{n_1>\cdots>n_d\ge1}
   \frac{x_1^{n_1}\cdots x_d^{n_d}}
        {n_1^{s_1}\cdots n_d^{s_d}},
 \qquad K_{6,1}=\operatorname{Im}\Li_{6,1}(i,-1).
\]
For precision about colors, write $x_j=i^{c_j}$, with $c_j$ read modulo
four. With $\omega_a=dt/(t-a)$ and outer-letter-first iterated integrals
$\mathcal I$, the corresponding word convention is
\begin{equation}
 \Li_{\boldsymbol s}(i^{c_1},\ldots,i^{c_d})
 =(-1)^d\mathcal I\left(
 0^{s_1-1}i^{-c_1}\cdots
 0^{s_d-1}i^{-(c_1+\cdots+c_d)}\right).
 \label{audit:s6-word-convention}
\end{equation}
All words below have first letter different from $1$ and last letter
different from $0$.

\begin{proposition}[A single-by-double shuffle reduction]
\label{audit:s6-depth-three}
The following identity consists entirely of convergent values:
\begin{align}
 K_{6,1}+\operatorname{Im}\Li_{5,2}(1,-i)
 ={}&\sum_{a=1}^{5}
       \operatorname{Im}\Li_{a,6-a,1}(i,-i,-1)\notag\\
 &+\operatorname{Im}\Li_{5,1,1}(1,i,i)
   +\operatorname{Im}\Li_{5,1,1}(1,-1,-i)\notag\\
 &-\operatorname{Im}\Li_{1,5,1}(i,1,-1)
   -\operatorname{Im}\Li_{5,1,1}(1,i,-1)\notag\\
 &-\operatorname{Im}\Li_{5,1,1}(1,-1,i).
 \label{audit:s6-depth-three-identity}
\end{align}
\end{proposition}
\begin{proof}
Consider $P=\Li_1(i)\Li_{5,1}(1,-1)$. Its stuffle expansion is
\begin{align*}
 P={}&\Li_{1,5,1}(i,1,-1)+\Li_{5,1,1}(1,i,-1)
       +\Li_{5,1,1}(1,-1,i)\\
    &+\Li_{6,1}(i,-1)+\Li_{5,2}(1,-i).
\end{align*}
For its integral shuffle, insert the single letter $-i$ into the word
$0^4(1)(-1)$ at each of its seven positions. Formula
\eqref{audit:s6-word-convention} then gives
\[
 P=\sum_{a=1}^{5}\Li_{a,6-a,1}(i,-i,-1)
       +\Li_{5,1,1}(1,i,i)+\Li_{5,1,1}(1,-1,-i).
\]
Subtract the two expansions and take imaginary parts. The first-index-one
terms have first argument $i$ and converge by summation by parts; all
remaining terms converge absolutely. Thus the finite-sum stuffle limit
and the convergent integral shuffle both apply. No divergent
$\Li_1(1)$ or regularization parameter is used.
\end{proof}

\paragraph{What this says about the known obstruction.}
The weight-seven functional $\ell$ in the pinned manuscript,
\texttt{cycloquot:eq:separating}, is a functional on a specified formal
depth-two presentation, not a functional on actual periods
\cite{Repo}. It annihilates the displayed single-by-single
shuffle/stuffle rows and the stated background; in particular,
\[
 \ell(K_{6,1})=1,\qquad
 \ell\bigl(\operatorname{Im}\Li_{5,2}(1,-i)\bigr)=0.
\]
Extend it to the free higher-depth word space by assigning zero to all
depth-three symbols, and call this particular extension $\widetilde\ell$.
For the convergent single-by-double row
\[
 R=\operatorname{Im}\{\mathsf{shuffle}(\Li_1(i),\Li_{5,1}(1,-1))
               -\mathsf{stuffle}(\Li_1(i),\Li_{5,1}(1,-1))\},
\]
the explicit expansions prove $\widetilde\ell(R)=-1$. Hence the naive
zero extension of the old separator already fails on this elementary
additional row. This calculation does not rule out another extension of
$\ell$, and it does not establish that any particular additional family,
including octahedral relations, is indispensable.

Let $Q_7$ denote the ten-term right side of
\eqref{audit:s6-depth-three-identity}. The ordinary stuffle of
$\Li_5(1)\Li_2(-i)$ gives
\[
 \operatorname{Im}\Li_{5,2}(1,-i)
 =g_{2,5}+\beta(7)-G\zeta(5),
 \qquad g_{a,b}=\operatorname{Im}\Li_{a,b}(i,1),\quad G=\beta(2).
\]
Selecting the even inner indices gives
$S_6=g_{6,1}+K_{6,1}$, so there is the exact reduction
\begin{equation}
 S_6=g_{6,1}-g_{2,5}+G\zeta(5)-\beta(7)+Q_7.
 \label{audit:s6-Q-reduction}
\end{equation}
The unresolved task is to reduce this specified $Q_7$ to the frozen
conjectural basket by a finite exact certificate. The double-shuffle
identities used here are classical; the contribution of this audit is
the explicit target row and the precise diagnostic of the inherited
obstruction.

The accompanying \texttt{verify\_s6\_target.py} independently enumerates
the seven integral shuffles and five stuffle terms using exact integer
arithmetic. It also checks all $192$ defining weight-seven restricted
rows against $\ell$, verifies the value $-1$ above, and records the
unchanged conjectural target in \texttt{s6\_target\_audit.json}.


\subsection{Three carry-forward manuscript corrections}

The following points were already identified in the incoming
Twisted Stieltjes report \cite{IncomingTwisted}. They remain relevant
at the pinned revision. The accompanying patch proposes only these
narrow changes and is supplied unapplied.

\paragraph{A real logarithm for the negative-parameter integral.}
In \texttt{chapters/03-algebraic.tex}, the displayed formula labelled
\texttt{cleo:eq:lintanh} uses $\log\tan(\theta/2)$ while its stated
real parameter range includes negative values. The real formula needs
$\log|\tan(\theta/2)|$. Under the principal complex logarithm,
the incorrect right side at $\lambda=-1/2$ differs from the real
integral by $-i\pi^2/6$. This is a branch correction, with the
absolute value fixed by the real integral and its real derivative.

\paragraph{Remove an unsupported nonreduction description.}
The same chapter calls a specific exact diagonal value
\emph{non-elementary} without defining a comparison algebra or proving
nonmembership. Keep the exact value and remove that classification.
An available special-function representation by itself proves no
nonreduction theorem.

\paragraph{Scope the coefficient field to the character in use.}
In \texttt{chapters/07-integration.tex}, the field
$\mathbb Q(\sqrt q)$ belongs to the particular quadratic-character
example. General additive or multiplicative character coordinates
can require cyclotomic and character values. The patch makes this
scope explicit without altering the quadratic calculation.

One earlier proposed wording change is already integrated: the PSLQ
discussion now allows known dependencies to be removed or quotiented
out and requires a nonzero coefficient of the target. It no longer
requires a prior linear-independence theorem. That edit is therefore
not repeated or counted as a new correction.

% The discussion concerns the inspected author PDF, not every published version.
% Bibliography keys to map in the parent: BorweinDilcher, DLMFpolylog.
\subsection{Corrections in the inspected Borwein--Dilcher author PDF}
\label{audit:BD}

The author-hosted PDF of Borwein and Dilcher,
\emph{Derivatives and fast evaluation of the Witten zeta function},
\url{https://carmamaths.org/resources/jon/omega3.pdf}, contains two
transcription errors relevant to a Mellin implementation.  Lemma~2.1(b),
printed page~2, assigns $H_0=1$; the required convention is $H_0=0$.
In Theorem~2.6, printed page~5, the positive-integer formula for $B_p$
has an extra factor $H_{p-1}$.  The corrected coefficients are
\begin{equation}\label{audit:BD-AB}
 A_p=\frac{(-1)^{p-1}H_{p-1}}{\Gamma(p)},\qquad
 B_p=\frac{(-1)^p}{\Gamma(p)},\qquad p=1,2,\ldots.
\end{equation}
These findings concern the inspected PDF; they do not assert that every
published version repeats them.

Here is an independent check.  Coalescing the Gamma pole and the
$\zeta(1)$ term in Jonqui\`ere's formula gives
\begin{equation}\label{audit:integer-Jonquiere}
 \operatorname{Li}_p(e^{-x})
 =\sum_{\substack{k\ge0\\k\ne p-1}}
       \zeta(p-k)\frac{(-x)^k}{k!}
 +\frac{(-x)^{p-1}}{(p-1)!}\bigl(H_{p-1}-\log x\bigr),
 \qquad 0<x<2\pi.
\end{equation}
Thus the coefficient multiplying $x^{p-1}\log x$ is precisely $B_p$
in \eqref{audit:BD-AB}.  At $p=3$, the required coefficient is $-1/2$,
whereas the printed $B_3$ is $-3/4$.  Separately,
$\operatorname{Li}_1(e^{-x})=-\log(1-e^{-x})=-\log x+O(x)$ forces
$H_0=0$.  A test at $p=2$ would miss the $B_p$ error because $H_1=1$.

Theorems~2.5 and~2.6 also state $\theta>0$ for their local power-series
expansions.  For ordinary convergent summation, impose
$0<\theta<2\pi$, as required by Jonqui\`ere's local series
\cite[25.12.12]{DLMFpolylog}.  A larger splitting point is possible only
with an additional continuation or intermediate-integral prescription.

To see the issue without an error estimate, take $r=s=0,t=3$ in the
double local series and keep $v=0,u=2k-1$.  The absolute value of this
summand is
\[
 \frac{|B_{2k}|\,\theta^{2k+2}}{2(2k)!(2k+2)}
 =\frac{\zeta(2k)\theta^2}{2k+2}
                       \left(\frac\theta{2\pi}\right)^{2k}.
\]
It fails to tend to zero for $\theta>2\pi$.  The original Witten sum at
$(0,0,3)$ converges, so the obstruction belongs to the displayed local
series representation.  The corrected implementation can retain the
classical theorem by choosing any $0<\theta<2\pi$; it does not need a
new assignment to divergent power series.


\subsection{Further transcription corrections in a Bailey--Borwein author PDF}
\label{audit:BB2014}

The following corrections concern the inspected author PDF of
\emph{Computation and theory of Mordell--Tornheim--Witten sums II},
dated 5 May 2014, printed pages 17--19 \cite{BB2014}. Equation numbers in
this subsection refer to that PDF. No assertion about other versions
is intended. Use its derivative convention
\begin{align}
 \omega_{a,b,c}(q,r,s)
 &=\sum_{n,m\ge1}
   \frac{(-\log n)^a(-\log m)^b(-\log(n+m))^c}
        {n^q m^r(n+m)^s},\notag\\
 \zeta_{a,b}(r,s)
 &=\sum_{k>j\ge1}\frac{(-\log k)^a(-\log j)^b}{k^rj^s}.
 \label{audit:BB-definitions}
\end{align}
Superscripts on $\Li$ and $\zeta$ below denote derivatives with respect
to their orders; $a,b,c$ are nonnegative integers.

\paragraph{Order-derivative signs in Eq.~(84).}
The factors printed as $(-1)^n\log^a n$ and $(-1)^m\log^b m$
must be $(-1)^a\log^a n$ and $(-1)^b\log^b m$.
Writing $\Gamma(s,z)=\int_z^\infty e^{-y}y^{s-1}\,dy$, the corrected
formula is
\begin{align}
 \omega_{a,b,c}(q,r,s)
 ={}&\sum_{n,m\ge1}\frac{(-\log n)^a(-\log m)^b}{n^q m^r}
     \partial_s^c\left\{
       \frac{\Gamma(s,n+m)}{\Gamma(s)(n+m)^s}\right\}\notag\\
 &+\int_{1/e}^{1}
    \partial_s^c\left\{\frac{(-\log u)^{s-1}}{\Gamma(s)}\right\}
    \Li_q^{(a)}(u)\Li_r^{(b)}(u)\,\frac{du}{u}.
 \label{audit:BB84}
\end{align}
Indeed $\partial_q^a n^{-q}=(-\log n)^a n^{-q}$.
Expanding the two polylogarithms in the integral from $0$ to $1/e$
and setting $y=-\log u$ gives the incomplete-Gamma term, including
the complete derivative of its Gamma ratio. There is no alternation
in $n$ or $m$.

\paragraph{The sign in Eq.~(86).}
The corrected identity is
\begin{equation}
 \omega_{0,0,b}(0,0,t)
 =\zeta^{(b)}(t-1)-\zeta^{(b)}(t).
 \label{audit:BB86}
\end{equation}
For $t>2$, exactly $k-1$ positive pairs have $n+m=k$, so its left
side is $\sum_{k\ge2}(k-1)(-\log k)^b/k^t$. This proves
\eqref{audit:BB86} directly. The printed right side has the opposite
sign; already at $b=0$ it is negative while the defining sum is positive.

\paragraph{Swapping orders also swaps derivative labels in Eqs.~(90)--(91).}
The corrected formulas are
\begin{align}
 \omega_{a,0,b}(s,0,t)+\omega_{b,0,a}(t,0,s)
 &=\zeta^{(a)}(s)\zeta^{(b)}(t)
        -\zeta^{(a+b)}(s+t),\label{audit:BB90}\\
 \omega_{a,b,0}(s,t,0)-\omega_{b,0,a}(t,0,s)
        -\omega_{a,0,b}(s,0,t)
 &=\zeta^{(a+b)}(s+t).\label{audit:BB91}
\end{align}
To see the indexing, $\omega_{a,0,b}(s,0,t)=\zeta_{b,a}(t,s)$.
Differentiate
$\zeta(t,s)+\zeta(s,t)=\zeta(s)\zeta(t)-\zeta(s+t)$
$a$ times in $s$ and $b$ times in $t$. The second omega term must
therefore carry the derivative labels $(b,0,a)$ when its spectral
arguments are $(t,0,s)$. The correction is invisible in the special
case $a=b$.

\paragraph{Gamma arguments in Eqs.~(93)--(94).}
Gamma functions and their derivatives are evaluated at the spectral
variable, not the number of differentiations. The corrected ladders are
\begin{align}
 \zeta_{b,a}(t,s)
 ={}&-\sum_{k=0}^{b-1}\binom bk
       \frac{\Gamma^{(b-k)}(t)}{\Gamma(t)}\zeta_{k,a}(t,s)\notag\\
 &+\frac1{\Gamma(t)}\int_0^1
   \frac{\Li_s^{(a)}(x)}{1-x}
   \log^b(-\log x)(-\log x)^{t-1}\,dx,
 \label{audit:BB93}\\
 \omega_{a,b,c}(q,r,s)
 ={}&-\sum_{k=0}^{c-1}\binom ck
       \frac{\Gamma^{(c-k)}(s)}{\Gamma(s)}\omega_{a,b,k}(q,r,s)\notag\\
 &+\frac1{\Gamma(s)}\int_0^1
   \Li_q^{(a)}(x)\Li_r^{(b)}(x)
   \log^c(-\log x)(-\log x)^{s-1}\,\frac{dx}{x}.
 \label{audit:BB94}
\end{align}
An empty sum is zero. For \eqref{audit:BB93}, differentiate
\[
 \Gamma(t)\zeta_{0,a}(t,s)
 =\int_0^1(-\log x)^{t-1}\frac{\Li_s^{(a)}(x)}{1-x}\,dx
\]
$b$ times with respect to $t$ and isolate the term with no derivative
on $\Gamma(t)$. The same Leibniz argument applied to
$\Gamma(s)\omega_{a,b,0}(q,r,s)$ gives \eqref{audit:BB94}.
Thus the occurrences of the Gamma argument $b$ in Eq.~(93) must be
$t$, and those of $c$ in Eq.~(94) must be $s$.

\paragraph{Direct convergence domains.}
These derivations require actual convergence before any continuation.
For real $q,r\ge0$, the integral formulas
\eqref{audit:BB84} and \eqref{audit:BB94} hold under
\begin{equation}
 s>0,\qquad q+s>1,\qquad r+s>1,\qquad q+r+s>2.
 \label{audit:BB-domain}
\end{equation}
The printed total-weight condition alone omits the two pair conditions:
at $(q,r,s)=(2,0,1)$ the $n=1$ subseries is harmonic and diverges.
Splitting the sum into $m\le n$ and $n<m$ proves sufficiency of the
three strict summability conditions; finite logarithmic powers do not
alter this open domain. Formula \eqref{audit:BB93} holds, for example,
when $t>1$ and $t+s>2$, while \eqref{audit:BB90}--\eqref{audit:BB91}
follow directly for $s,t>1$. The latter identities and
\eqref{audit:BB86} then extend meromorphically. Such continuation does
not assert ordinary convergence of the displayed integrals or series
outside their stated domains.


\subsection{Conventions that must survive integration}

No new false theorem was found among the inspected periodic Hurwitz,
coordinate-residue, and coincident-pair premises used in this article.
The following assumptions are nevertheless part of the mathematics
and must remain attached to the statements.

\begin{enumerate}
\item The pointwise product finite part, the periodic distributional
      convolution, and the canonical extension in the shift coordinate
      are different constructions. Theorem~\ref{cont:main} computes
      their relation, including its contact term.
\item The argument finite part uses a specified cutoff coordinate.
      In Theorem~\ref{cub:general-completion}, subtracting only the
      residue restricted to a polar hyperplane loses finite numerator
      derivatives and changes the requested answer.
\item A fixed-integer identity may be differentiated only in variables
      that remain free in that identity. In particular, the finite
      positive-integer Tornheim sums in \eqref{sp:positive} do not
      determine derivatives in either fixed inner order.
\item An origin ray is restricted before its removable singularity is
      filled in. The two distinct values $1/3$ and $5/12$ are both
      correct for their respective rays.
\item Numerical verification checks an implementation and can detect
      mistakes. It does not replace the proofs or establish arithmetic
      independence.
\end{enumerate}

The source audit and correction patch are separate from the new
mathematical section files. Their provenance notes identify which
findings are inherited, which are newly checked against an external
PDF, and which are restrictions on interpretation.
